\chapter{Discussion}
\label{ch:discussion}

\section{What the degradation architecture bought, and what it did not}

The premise in \cref{ch:intro} was that a body-worn master will lose channels,
so losing them must be survivable. The two dated captures support the premise
empirically: the failed population grew between them
(\cref{sec:channelhealth}). The architecture's value can be stated precisely
from the measurements.

It kept the system running. With eight of fourteen channels failing the health
criterion, the right arm retains full position sensing at $R^2 = \num{0.986}$
and the left arm retains direction. A joint-space mapping would have been
inoperable on either arm from the first failed channel, because it has no
reconstruction path.

It kept the failures visible. The incoherent verdict exists because a failing
potentiometer does not go quiet, and range alone calls a wide-spanning
nonsense signal healthy. Without that verdict the system would have kept
following \SI{54.5}{\milli\metre} of commanded motion the operator never
asked for, on a master whose whole reach is \SI{272}{\milli\metre}.

What it did not buy is capability. Freezing the left arm's bend joints removes
its radial dimension entirely, collapsing its commanded set from a shell to a
surface, and the wrist-roll fallback restores an axis rather than restoring
the measurement. The architecture converts a fault into a known, bounded,
announced reduction. It does not convert it into working hardware, and
\cref{sec:channelhealth} names the two channels whose repair would.

\section{The reduced-degree-of-freedom premise, assessed}

The second premise was that commanding fourteen joints through a body-worn
master is a usability problem, and that the answer is to take each component
from the sensor that observes it best. The evidence is mixed, and it divides
cleanly by component.

Elevation from gravity works. It is mount-independent, drift-free, integrates
nothing, and survived every reformulation. Reach from the forward-kinematic
magnitude works, because a magnitude depends on the bend joints and tolerates
a dead roll passed through as zero. Both vertical and fore-and-aft motion
track correctly on both arms.

Azimuth from a single joint does not work, and the reason is instructive
rather than incidental. The spherical model assumes azimuth and elevation are
independent. They are not, because joint one is a roll about the arm's own
axis, so its effect on the tip scales with how bent the arm is. The
\SI{137.7}{\degree} measured between two directions intended to be
\SI{90}{\degree} apart is not a calibration error that a better axis map or a
better zero could remove; it is the model being wrong about the geometry.

The honest reading is that the decomposition succeeded for the two components
where a single sensor observes the quantity directly, and failed for the one
where it does not. That is a usable design rule for a successor: decompose
along the axes your sensors actually observe, not along the axes that make the
mathematics tidy.

\section{The workspace result and what it costs the study}

\Cref{ch:workspace} reports that the two arms share no reachable workspace
under the as-built mount. The consequence for the experimental programme is
severe and worth stating plainly rather than absorbing into a limitations
list.

Inter-arm handover is impossible. Not difficult, not unreliable: there is no
point in space both arms can reach, so no control strategy, no planner and no
orientation mode can produce one. The task that would have demonstrated
autonomy achieving something teleoperation cannot was blocked by geometry
before it was blocked by anything else, and the finding is the reason it is
not in the task set.

The cause is not the mount rotation, which was proved to leave the mirror
residual invariant. It is that the two arms are parked asymmetrically at their
home joint angles, and those angles are ground truth because the simulation to
hardware bridge replays them onto the real arm. Changing them in software
would command the difference as a jump. Fixing this requires re-parking the
right arm in hardware.

Two of the five delivered tasks are the same bimanual mechanism with different
objects, and only one task uses the second mechanism. There are two genuinely
bimanual mechanisms available on this platform, not three. That is a real
narrowing of what the platform can be used to ask.

\section{Comparison against the planning report's targets}

The delivered system diverges substantially from the plan, and the divergences
fall into three groups.

\emph{Superseded by measurement.} The planning report specified physical
springs at the master for force feedback and virtual spring impedance at the
slave. The impedance target is unreachable on this platform for a reason
discovered during the work: the arms' cyclic control interface is unusable
over the available network path, so the system commands through a high-level
velocity interface at \SIrange{18.4}{18.7}{\hertz}. An impedance law at
\SI{18}{\hertz} is not an impedance law. The compliance the plan sought would
have to come from the mechanism instead, through series elastic actuation at
the joint~\cite{toubar2022sea,lee2019sea}, or from a current-based
formulation that does not require a high-rate torque
loop~\cite{dewolde2024impedance}. Both are hardware changes.

\emph{Deferred and still open.} Coordinated pick-and-place on the real
platform was not demonstrated. The task set is specified and verified
geometrically, and the runner exists, but the arms were not available for the
period in which it would have been run.

\emph{Added.} The safety architecture, the degradation architecture and the
workspace characterisation were not in the plan. The first two are the direct
consequence of discovering that sensor faults dominate; the third is the
consequence of asking whether the task set was reachable before running it.

\section{Threats to validity}

\subsection{Simulation is not the arm}

Every quantitative result in this thesis except the channel measurements comes
from simulation with mock hardware. The mock echoes commanded positions with
no dynamics, which has one consequence that must not be glossed: every force
reading is identically zero, so any quantity derived from a force is
structurally undefined rather than approximated. The measurement that payload
configuration made exactly \SI{0}{\metre} of difference to tracking is a
statement about the mock.

\subsection{Latency figures are lower bounds}

Every latency reported is control-path latency, measured between software
stages. It excludes the operator's reaction, the display, and the physical
arm's response. Any end-to-end figure will be larger, and the reported values
should be read as a floor.

\subsection{The instrument failure rate}

Four times, a result that looked like a system failure was a defect in the
measuring tool, and three of the four were caught only by checking the
instrument against a known answer. The base rate is high enough that any
unreproduced single measurement in this thesis should be treated as
provisional. This is the reason the verification methodology exists and the
reason the task-set figures are reported with their repeat count.

\subsection{The one-stack constraint}

Two instances of the master node split the serial stream between them and
invalidated a full day of measurements before the cause was found. This is now
prevented at the file-descriptor level, but it means that any measurement
taken before that prevention was added carries a risk that cannot be
retrospectively excluded.

\section{What stands between this system and a two-person study}

The study designed in \cref{ch:method} cannot be run today, and the blockers
are specific.

\begin{enumerate}
  \item \textbf{Two master channels.} Left j2 and left j4. With both frozen
    the left arm has no radial dimension, so the direct teleoperation
    condition, which is the baseline the whole comparison rests on, is not
    representative of a working master.
  \item \textbf{The right arm's home pose.} It has never been read from
    hardware, and the workspace geometry depends on it.
  \item \textbf{Hardware validation of the safety layer.} Nothing in the
    safety architecture has run against a real arm. The Kortex session
    recovery path is the riskiest part, because the arm permits exactly one
    session and a leaked one refuses the next connection.
  \item \textbf{Ethics approval for a two-person protocol}, which differs from
    a single-participant protocol in that one participant bears a physical
    risk arising from another participant's actions.
    \todo{CONFIRM: current status of the ethics application with Dr.\ Zhang.}
  \item \textbf{A detection rate measured on real images.} The synthetic
    figure characterises the renderer.
\end{enumerate}

None of these is a research question. All are execution, and the first is two
soldered joints.
